% secciones/06_conclusiones.tex
\section{Conclusiones y Trabajo Futuro}
La realización de este proyecto ha permitido abordar de forma integral el desafío que supone el diseño de un sistema robótico autónomo para un entorno tan crítico y exigente como es el hospitalario. La integración de sistemas robóticos móviles en la logística hospitalaria representa un cambio de paradigma en la gestión de servicios asistenciales, permitiendo una optimización sin precedentes de los flujos de materiales críticos. A lo largo de esta investigación, se ha evidenciado que el éxito de estos despliegues no depende únicamente de la sofisticación técnica del hardware, sino de una visión sistémica que contemple la seguridad biológica, el cumplimiento estricto de la normativa ISO 13482 y una integración profunda con la infraestructura inteligente del centro sanitario. La adopción de ROS 2 como núcleo tecnológico facilita la creación de sistemas modulares y escalables, capaces de adaptarse a las necesidades dinámicas de un hospital moderno.

Una de las principales conclusiones técnicas es la idoneidad de las ruedas Mecanum para la logística interna. La capacidad de traslación lateral no es solo una ventaja de maniobrabilidad, sino un factor de seguridad que permite al robot apartarse rápidamente en pasillos congestionados ante el paso de una camilla de urgencias. Asimismo, la arquitectura de computación híbrida (Jetson Orin + STM32) garantiza que el sistema sea capaz de procesar algoritmos complejos de visión artificial sin comprometer el determinismo necesario para el control de los motores y la seguridad física.

Desde una perspectiva operativa, la implementación de estaciones de desinfección automatizada y el cumplimiento de normativas como la ISO 13482 y el estándar IEC 62443 son los pilares que transforman un prototipo de laboratorio en una herramienta profesional confiable. La automatización de la logística no debe verse como un fin en sí mismo, sino como un medio para revalorizar el trabajo humano, permitiendo que el personal de enfermería recupere su rol asistencial y humano, delegando las tareas de carga y transporte en sistemas autónomos seguros y trazables.

\subsection{Trabajo Futuro}
A pesar de la robustez del diseño propuesto, existen diversas líneas de investigación y desarrollo que podrían ampliar las capacidades del sistema en el futuro:
\begin{itemize}
    \item \textbf{Coordinación de Flotas Multirobot:} Implementar algoritmos de gestión de tráfico centralizados mediante \textit{Open-RMF} para coordinar una flota de decenas de unidades, optimizando el uso de ascensores y evitando cuellos de botella en las zonas de carga.
    \item \textbf{Integración con el Historial Clínico:} Vincular el sistema de trazabilidad del robot con el software de gestión hospitalaria (HIS) para automatizar la solicitud de transporte de muestras basándose en las órdenes médicas en tiempo real.
    \item \textbf{Mejora de la Interacción Social:} Desarrollar modelos de lenguaje de gran tamaño (LLM) optimizados para su ejecución en el borde (\textit{Edge}) que permitan al robot interactuar de forma natural con pacientes y visitantes, resolviendo dudas simples o proporcionando guiado por voz.
    \item \textbf{Optimización del Consumo Energético:} Investigar el uso de algoritmos de navegación que tengan en cuenta el coste energético de las maniobras, priorizando rutas que maximicen la vida útil de la batería en función de la urgencia de la misión.
\end{itemize}

En definitiva, este trabajo sienta las bases para un despliegue realista de la robótica de servicio en la sanidad pública, demostrando que la tecnología actual está madura para asumir roles logísticos críticos, siempre que se diseñen bajo principios de seguridad, asepsia y respeto a la privacidad del paciente.
